%==============================================================================
%  CHAPTER 12: Advanced Optimizers
%  Part II: Optimization Theory
%==============================================================================
\documentclass[../../main.tex]{subfiles}

\begin{document}

\chapter{Advanced Optimizers}
\label{ch:advanced-optimizers}

Adam is used in the vast majority of deep learning papers published today. But in 2018, researchers discovered that Adam can fail catastrophically on certain problems---while plain SGD succeeds (Reddi et al., 2018). Knowing when to use which optimizer is the difference between state-of-the-art results and wasted compute.

\begin{whycare}
Modern optimizers like Adam\index{Adam}, AdamW\index{AdamW}, and LAMB\index{LAMB|textbf} encode decades of research into a few lines of code. Understanding their internals lets you diagnose training failures, tune hyperparameters intelligently, and choose the right tool for each job.
\end{whycare}

\begin{keyconcept}
This chapter is essential reading for practitioners. It covers the adaptive optimizers (AdaGrad, RMSprop, Adam, AdamW) that power modern deep learning, along with learning rate schedules, optimizer selection, hyperparameter tuning, and failure mode diagnosis. Building on the SGD and momentum foundations from Chapter~\ref{ch:gradient-descent}, we explore both the theory of \textit{why} these algorithms work and the practical guidance needed to apply them effectively.
\end{keyconcept}

\begin{realworld}
\textbf{Why This Matters:} Adam is the default optimizer for good reason---it adapts learning rates per-parameter automatically.

\textbf{Real Example:} Adam has become the default optimizer for most deep learning research, with variants like AdamW and LAMB widely adopted.

\textbf{Scale:} AdamW with learning rate warmup enabled stable training of the largest language models.
\end{realworld}

\begin{prerequisites}
\textbf{Before reading this chapter, you should be comfortable with:}
\begin{itemize}
    \item Gradient descent variants: SGD, momentum (Chapter~\ref{ch:gradient-descent})
    \item Exponential moving averages
    \item Basic convergence concepts (Chapter~\ref{ch:stochastic-optimization})
\end{itemize}
\textbf{Review if needed:} Chapter~\ref{ch:gradient-descent} for SGD and momentum basics
\end{prerequisites}

%------------------------------------------------------------------------------
\section{Introduction}
\label{sec:adv-opt-intro}

Building on the gradient descent and momentum methods covered in Chapter~\ref{ch:gradient-descent}, this chapter explores adaptive learning rate methods that automatically adjust per-parameter learning rates based on historical gradient information.

While momentum methods (see Section~\ref{sec:ch11-momentum} in Chapter~\ref{ch:gradient-descent}) accelerate convergence by accumulating gradient information over time, they still use a single global learning rate for all parameters. Different parameters in a neural network often benefit from different learning rates---parameters associated with frequently occurring features might need smaller learning rates, while parameters for rare features need larger rates. Adaptive methods address this limitation by automatically scaling the learning rate for each parameter.

This chapter covers AdaGrad, RMSprop, Adam, and AdamW---the adaptive optimizers that power most modern deep learning. We also explore learning rate scheduling, optimizer selection guidelines, and practical hyperparameter tuning strategies.

%------------------------------------------------------------------------------
\section{Adaptive Learning Rate Methods}
\label{sec:adaptive-lr}

Different parameters in a neural network may benefit from different learning rates. Parameters associated with frequently occurring features might need smaller learning rates (they update often), while parameters for rare features need larger rates (they update infrequently). \textbf{Adaptive methods} automatically adjust per-parameter learning rates based on historical gradient information.

\begin{keyconcept}
\textbf{The Adaptive Optimizer Family Tree:}

The adaptive optimizers form a natural lineage, where each algorithm addressed limitations of its predecessor:

\begin{center}
\small
\begin{tikzpicture}[node distance=1.2cm, auto,
    box/.style={rectangle, draw, rounded corners, minimum width=2cm, minimum height=0.7cm, align=center, font=\footnotesize}]
    \node[box] (adagrad) {AdaGrad\\(2011)};
    \node[box, right=of adagrad] (rmsprop) {RMSprop\\(2012)};
    \node[box, right=of rmsprop] (adam) {Adam\\(2014)};
    \node[box, right=of adam] (adamw) {AdamW\\(2017)};
    \draw[->, thick] (adagrad) -- (rmsprop);
    \draw[->, thick] (rmsprop) -- (adam);
    \draw[->, thick] (adam) -- (adamw);
\end{tikzpicture}
\end{center}

\textbf{The evolution:}
\begin{itemize}
    \item \textbf{AdaGrad} (Duchi et al., 2011): First to adapt learning rates per-parameter based on accumulated squared gradients. \textit{Problem}: Learning rate decays to zero over time, stopping learning.
    \item \textbf{RMSprop} (Hinton, 2012): Uses exponential moving average instead of sum, so learning rate doesn't vanish. \textit{Problem}: No momentum component.
    \item \textbf{Adam} (Kingma \& Ba, 2014): Combines RMSprop's adaptive rates with momentum, plus bias correction. \textit{Problem}: Weight decay interacts poorly with adaptive rates.
    \item \textbf{AdamW} (Loshchilov \& Hutter, 2017): Decouples weight decay from gradient updates, fixing generalization issues.
\end{itemize}

Each step solved a real problem practitioners encountered. Understanding this history helps you diagnose when to deviate from defaults.
\end{keyconcept}

\subsection{AdaGrad}

\begin{definition}{AdaGrad}{adagrad-def}
\textbf{Adaptive Gradient}\index{AdaGrad|textbf} (AdaGrad)~\cite{duchi2011adaptive} accumulates squared gradients to scale the learning rate for each parameter:
\begin{align}
    \vect{G}_{t+1} &= \vect{G}_t + \vect{g}_t \odot \vect{g}_t \\
    \vect{\theta}_{t+1} &= \vect{\theta}_t - \frac{\eta}{\sqrt{\vect{G}_{t+1}} + \epsilon} \odot \vect{g}_t
\end{align}
where $\vect{g}_t = \grad L(\vect{\theta}_t)$, $\odot$ denotes element-wise multiplication, and $\epsilon \approx 10^{-8}$ prevents division by zero.
\end{definition}

AdaGrad's key insight is that parameters with large accumulated gradients receive smaller effective learning rates, while those with small accumulated gradients receive larger rates. This is particularly beneficial for:
\begin{itemize}
    \item \textbf{Sparse features}: Common in NLP where some words appear frequently while others are rare
    \item \textbf{Ill-conditioned problems}: Different parameters may have vastly different optimal scales
\end{itemize}

\begin{important}
\textbf{AdaGrad's Limitation:} The accumulated sum $\vect{G}_t$ grows monotonically, causing the effective learning rate to shrink continually. For long training runs, the learning rate may become vanishingly small before the model converges, essentially stopping learning. This motivates RMSprop and Adam.
\end{important}

\begin{algorithm}[htbp]
\caption{AdaGrad}
\label{alg:adagrad}
\begin{algorithmic}[1]
\Require Loss function $L$, initial $\vect{\theta}_0$, learning rate $\eta$, $\epsilon = 10^{-8}$
\Ensure Optimized parameters $\vect{\theta}^*$
\State Initialize $\vect{G}_0 \gets \mathbf{0}$
\For{$t = 0, 1, 2, \ldots$ until convergence}
    \State Compute gradient: $\vect{g}_t \gets \grad L(\vect{\theta}_t)$
    \State Accumulate squared gradients: $\vect{G}_{t+1} \gets \vect{G}_t + \vect{g}_t \odot \vect{g}_t$
    \State Update: $\vect{\theta}_{t+1} \gets \vect{\theta}_t - \frac{\eta}{\sqrt{\vect{G}_{t+1}} + \epsilon} \odot \vect{g}_t$
\EndFor
\State \Return $\vect{\theta}_t$
\end{algorithmic}
\end{algorithm}

\subsection{RMSprop}

\begin{definition}{RMSprop}{rmsprop-def}
\textbf{Root Mean Square Propagation}\index{RMSprop|textbf} (RMSprop), proposed by Hinton in lecture slides for his 2012 Coursera course, addresses AdaGrad's diminishing learning rate by using an exponentially decaying average of squared gradients:
\begin{align}
    \vect{v}_t &= \beta \vect{v}_{t-1} + (1 - \beta) \vect{g}_t \odot \vect{g}_t \\
    \vect{\theta}_{t+1} &= \vect{\theta}_t - \frac{\eta}{\sqrt{\vect{v}_t} + \epsilon} \odot \vect{g}_t
\end{align}
where $\beta$ (called \texttt{alpha} in PyTorch) defaults to $0.99$ in PyTorch, though $0.9$ is also common.
\end{definition}

The exponential moving average ensures that $\vect{v}_t$ represents the \textit{recent} magnitude of gradients rather than their entire history. This prevents the learning rate from decaying to zero while still providing adaptive scaling.

\begin{intuition}
RMSprop can be understood as keeping a running estimate of gradient variance. Parameters with consistently large gradients get smaller learning rates (they are already moving fast), while parameters with small gradients get larger rates (they need more encouragement to move).
\end{intuition}

\begin{commonmistakes}
\begin{itemize}
    \item \textbf{Mistake:} Confusing RMSprop's $\beta$ with momentum's $\gamma$ (Chapter~\ref{ch:gradient-descent})

    \textbf{Why it's wrong:} Both use exponential moving averages, but they track different quantities. Momentum's $\gamma$ smooths the gradient \textit{direction}, while RMSprop's $\beta$ smooths the gradient \textit{magnitude squared}. These serve fundamentally different purposes.

    \textbf{Correct understanding:} Momentum accelerates convergence along consistent directions. RMSprop adapts the step size per-parameter based on historical gradient magnitudes. Adam combines both ideas.
\end{itemize}
\end{commonmistakes}

\begin{algorithm}[htbp]
\caption{RMSprop}
\label{alg:rmsprop-ch12}
\begin{algorithmic}[1]
\Require Loss function $L$, initial $\vect{\theta}_0$, learning rate $\eta = 0.01$, decay $\beta = 0.99$, $\epsilon = 10^{-8}$
\Ensure Optimized parameters $\vect{\theta}^*$
\State Initialize $\vect{v}_0 \gets \mathbf{0}$
\For{$t = 0, 1, 2, \ldots$ until convergence}
    \State Compute gradient: $\vect{g}_t \gets \grad L(\vect{\theta}_t)$
    \State Update squared gradient average: $\vect{v}_t \gets \beta \vect{v}_{t-1} + (1-\beta) \vect{g}_t \odot \vect{g}_t$
    \State Update: $\vect{\theta}_{t+1} \gets \vect{\theta}_t - \frac{\eta}{\sqrt{\vect{v}_t} + \epsilon} \odot \vect{g}_t$
\EndFor
\State \Return $\vect{\theta}_t$
\end{algorithmic}
\end{algorithm}

\subsection{Exponential Moving Averages}\label{sec:ema-intro}

Before diving into adaptive optimizers, we need one key concept: the \textbf{exponential moving average} (EMA).

\begin{definition}{Exponential Moving Average}{ema}
Given a sequence of values $x_1, x_2, \ldots$, the exponential moving average with decay $\beta$ is:
\[
v_t = \beta \cdot v_{t-1} + (1 - \beta) \cdot x_t
\]
Starting from $v_0 = 0$.
\end{definition}

\begin{intuition}{EMA as Recency-Weighted Average}
The EMA is a running average that weights recent values more heavily. The parameter $\beta$ controls the ``memory'':
\begin{itemize}
    \item $\beta = 0.9$: Roughly averages the last $\frac{1}{1-\beta} = 10$ values
    \item $\beta = 0.99$: Roughly averages the last 100 values
    \item $\beta = 0.999$: Roughly averages the last 1000 values
\end{itemize}
Higher $\beta$ means smoother estimates but slower adaptation to changes.
\end{intuition}

\textbf{Example}: With $\beta = 0.9$ and sequence $[1, 2, 3, 4, 5]$:
\begin{align*}
v_1 &= 0.9 \cdot 0 + 0.1 \cdot 1 = 0.1 \\
v_2 &= 0.9 \cdot 0.1 + 0.1 \cdot 2 = 0.29 \\
v_3 &= 0.9 \cdot 0.29 + 0.1 \cdot 3 = 0.561 \\
&\vdots
\end{align*}

Adam uses \textit{two} EMAs: one for the gradient (first moment, $m$) and one for the squared gradient (second moment, $v$). This is why understanding EMAs is essential.

\subsection{Adam: Adaptive Moment Estimation}

\begin{definition}{Adam Optimizer}{adam-def}
\textbf{Adam}\index{Adam|textbf}\index{adaptive moment estimation|see{Adam}}~\cite{kingma2015adam} combines momentum with adaptive learning rates by tracking both first and second moments of gradients:
\begin{align}
    \vect{m}_t &= \beta_1 \vect{m}_{t-1} + (1 - \beta_1) \vect{g}_t & \text{(first moment)}\index{first moment|textbf} \\
    \vect{v}_t &= \beta_2 \vect{v}_{t-1} + (1 - \beta_2) \vect{g}_t \odot \vect{g}_t & \text{(second moment)}\index{second moment|textbf} \\
    \hat{\vect{m}}_t &= \frac{\vect{m}_t}{1 - \beta_1^t} & \text{(bias-corrected first moment)} \\
    \hat{\vect{v}}_t &= \frac{\vect{v}_t}{1 - \beta_2^t} & \text{(bias-corrected second moment)} \\
    \vect{\theta}_{t+1} &= \vect{\theta}_t - \frac{\eta}{\sqrt{\hat{\vect{v}}_t} + \epsilon} \odot \hat{\vect{m}}_t
\end{align}
Default hyperparameters: $\beta_1 = 0.9$, $\beta_2 = 0.999$, $\eta = 0.001$, $\epsilon = 10^{-8}$\index{epsilon|textbf}.
\end{definition}

\begin{keyconcept}
Adam's bias correction\index{bias correction|textbf} is crucial for stable early training. Since $\vect{m}_0 = \vect{v}_0 = \mathbf{0}$, the initial moment estimates are biased toward zero. The correction factors $1/(1 - \beta^t)$ compensate for this bias. Early in training when $t$ is small, $\beta^t$ is close to 1, making the correction significant. As training progresses, $\beta^t \to 0$ and the correction becomes negligible.
\end{keyconcept}

\begin{aha}
Adam = Momentum + RMSprop + bias correction. Momentum remembers which direction you've been going. RMSprop adapts step size per parameter. Bias correction fixes the cold-start problem in early iterations. Three ideas, one powerful optimizer.
\end{aha}

\begin{checkpoint}
Why does Adam need bias correction? What would happen to $m_1$ and $v_1$ without it?
\end{checkpoint}

\begin{algorithmtrace}
\textbf{Example: Adam on} $f(x) = x^2$, $x_0 = 4$, $\eta=0.5$, $\beta_1=0.9$, $\beta_2=0.999$

\textbf{Iteration 1:} $g = 2x = 8.0$
  $m_1 = 0.9 \times 0 + 0.1 \times 8.0 = 0.8$
  $v_1 = 0.999 \times 0 + 0.001 \times 64 = 0.064$
  $\hat{m}_1 = 0.8 / (1 - 0.9^1) = 8.0$ \quad (bias correction!)
  $\hat{v}_1 = 0.064 / (1 - 0.999^1) = 64.0$
  $x_1 = 4.0 - 0.5 \times 8.0 / (\sqrt{64.0} + 10^{-8}) = 3.5$

\textbf{Iteration 2:} $g = 2x = 7.0$
  $m_2 = 0.9 \times 0.8 + 0.1 \times 7.0 = 1.42$
  $v_2 = 0.999 \times 0.064 + 0.001 \times 49 = 0.113$
  $\hat{m}_2 = 1.42 / (1 - 0.81) = 7.47$
  $\hat{v}_2 = 0.113 / (1 - 0.998) = 56.5$
  $x_2 = 3.5 - 0.5 \times 7.47 / (\sqrt{56.5} + 10^{-8}) = 3.00$

Note: Bias correction is crucial in early iterations!
Without it, $m_1 = 0.8$ would underestimate the true gradient of 8.0.
\end{algorithmtrace}

To understand why bias correction is needed, consider the expected value of $\vect{m}_t$ assuming stationary gradients with mean $\vect{\mu}$:
\[
    \E[\vect{m}_t] = (1 - \beta_1^t) \vect{\mu}
\]
The bias-corrected estimate $\hat{\vect{m}}_t = \vect{m}_t / (1 - \beta_1^t)$ has expected value $\vect{\mu}$, restoring unbiasedness.

\begin{algorithm}[htbp]
\caption{Adam Optimizer}
\label{alg:adam-ch12}
\begin{algorithmic}[1]
\Require Loss function $L$, initial $\vect{\theta}_0$
\Require $\eta = 0.001$, $\beta_1 = 0.9$, $\beta_2 = 0.999$, $\epsilon = 10^{-8}$
\Ensure Optimized parameters $\vect{\theta}^*$
\State Initialize $\vect{m}_0 \gets \mathbf{0}$, $\vect{v}_0 \gets \mathbf{0}$
\For{$t = 1, 2, \ldots$ until convergence}
    \State Compute gradient: $\vect{g}_t \gets \grad L(\vect{\theta}_{t-1})$
    \State Update first moment: $\vect{m}_t \gets \beta_1 \vect{m}_{t-1} + (1-\beta_1) \vect{g}_t$
    \State Update second moment: $\vect{v}_t \gets \beta_2 \vect{v}_{t-1} + (1-\beta_2) \vect{g}_t \odot \vect{g}_t$
    \State Bias correction: $\hat{\vect{m}}_t \gets \vect{m}_t / (1 - \beta_1^t)$
    \State Bias correction: $\hat{\vect{v}}_t \gets \vect{v}_t / (1 - \beta_2^t)$
    \State Update: $\vect{\theta}_t \gets \vect{\theta}_{t-1} - \eta \cdot \hat{\vect{m}}_t / (\sqrt{\hat{\vect{v}}_t} + \epsilon)$
\EndFor
\State \Return $\vect{\theta}_t$
\end{algorithmic}
\end{algorithm}

\begin{pythoncode}
import numpy as np

class Adam:
    """Adam optimizer implementation."""
    def __init__(self, lr=0.001, beta1=0.9, beta2=0.999, eps=1e-8):
        self.lr, self.beta1, self.beta2, self.eps = lr, beta1, beta2, eps
        self.m, self.v, self.t = None, None, 0

    def update(self, params, grads):
        """Perform one Adam update step."""
        if self.m is None:
            self.m, self.v = np.zeros_like(params), np.zeros_like(params)
        self.t += 1
        self.m = self.beta1 * self.m + (1 - self.beta1) * grads
        self.v = self.beta2 * self.v + (1 - self.beta2) * grads**2
        # Bias correction compensates for initialization at zero.
        # At t=1 with beta1=0.9: m_1 = 0.1*g (biased low)
        # Correction: m_hat = m_1/(1-0.9^1) = m_1/0.1 = g (unbiased)
        m_hat = self.m / (1 - self.beta1**self.t)
        v_hat = self.v / (1 - self.beta2**self.t)
        return params - self.lr * m_hat / (np.sqrt(v_hat) + self.eps)
\end{pythoncode}

\begin{theorem}{Adam Update Bound}{adam-bound}
Under the assumption that gradients are bounded $\norm{g_t}_\infty \leq G$ for some constant $G$, Adam's effective step size is bounded:
\[
    \left|\frac{\eta \hat{m}_{t,i}}{\sqrt{\hat{v}_{t,i}} + \epsilon}\right| \leq \frac{\eta}{\sqrt{1-\beta_2}}
\]
This bound (from Kingma \& Ba, 2015) ensures that individual parameter updates remain controlled regardless of gradient scale, providing implicit gradient clipping behavior.
\end{theorem}

\subsection{AdamW: Decoupled Weight Decay}

\begin{definition}{AdamW}{adamw-def}
\textbf{AdamW}\index{AdamW|textbf}\index{decoupled weight decay|textbf} (Loshchilov \& Hutter, 2019) decouples weight decay regularization from the adaptive gradient update:
\[
    \vect{\theta}_{t+1} = \vect{\theta}_t - \eta \left( \frac{\hat{\vect{m}}_t}{\sqrt{\hat{\vect{v}}_t} + \epsilon} + \lambda \vect{\theta}_t \right)
\]
where $\lambda$ is the weight decay coefficient. This differs from L2 regularization in standard Adam.
\end{definition}

\begin{important}
\textbf{Weight Decay vs. L2 Regularization:} In standard SGD, adding L2 regularization $\frac{\lambda}{2}\norm{\vect{\theta}}^2$ to the loss is equivalent to adding $\lambda \vect{\theta}$ to the gradient (weight decay). However, in adaptive optimizers like Adam, these are \textit{not} equivalent. L2 regularization adds $\lambda \vect{\theta}$ to $\vect{g}_t$ before the adaptive scaling, while decoupled weight decay applies it after. AdamW uses decoupled weight decay, which empirically leads to better generalization.

\textbf{Why Decoupling Matters:} With L2 regularization in Adam, the regularization gradient $\lambda \vect{\theta}$ gets scaled by the adaptive learning rate $1/\sqrt{\hat{\vect{v}}_t}$. This means parameters with historically large gradients (large $\hat{\vect{v}}_t$) receive \emph{weaker} regularization, while parameters with small gradients receive \emph{stronger} regularization---the opposite of what we typically want. Decoupled weight decay in AdamW applies the penalty $\lambda \vect{\theta}_t$ directly to the weights without adaptive scaling, ensuring uniform regularization strength across all parameters regardless of their gradient history.
\end{important}

\begin{checkpoint}
What's the difference between Adam and AdamW? When does it matter?
\end{checkpoint}

The distinction matters because:
\begin{itemize}
    \item With L2 regularization in Adam, the regularization term is scaled by $1/\sqrt{\hat{v}_t}$, weakening regularization for parameters with large gradients
    \item With decoupled weight decay, all parameters experience the same relative regularization strength
\end{itemize}

\begin{pythoncode}
# These behave DIFFERENTLY despite similar syntax:
# Adam with L2 regularization (couples weight decay with adaptive LR)
opt_adam = torch.optim.Adam(model.parameters(), lr=0.001, weight_decay=0.01)

# AdamW with decoupled weight decay (recommended for transformers)
opt_adamw = torch.optim.AdamW(model.parameters(), lr=0.001, weight_decay=0.01)
# AdamW applies weight decay directly to weights, not through gradients
\end{pythoncode}

\begin{pythoncode}
class AdamW(Adam):
    """AdamW optimizer with decoupled weight decay."""
    def __init__(self, lr=0.001, beta1=0.9, beta2=0.999, eps=1e-8, wd=0.01):
        super().__init__(lr, beta1, beta2, eps)
        self.wd = wd

    def update(self, params, grads):
        """Perform one AdamW update step with decoupled weight decay."""
        if self.m is None:
            self.m, self.v = np.zeros_like(params), np.zeros_like(params)
        self.t += 1
        self.m = self.beta1 * self.m + (1 - self.beta1) * grads
        self.v = self.beta2 * self.v + (1 - self.beta2) * grads**2
        m_hat = self.m / (1 - self.beta1**self.t)
        v_hat = self.v / (1 - self.beta2**self.t)
        return params - self.lr * (m_hat / (np.sqrt(v_hat) + self.eps) + self.wd * params)
\end{pythoncode}

\begin{warstory}
\textbf{When Adam Failed}

When researchers trained BERT-scale models with typical batch sizes (around 256), Adam worked well. When they scaled to very large batch sizes (tens of thousands of samples) to dramatically reduce training time, Adam often failed. Gradients became unstable, and models wouldn't converge.

The problem? Adam's adaptive learning rates are calibrated for small-batch gradient statistics. At massive batch sizes, the gradient estimates are too accurate---there's not enough noise for Adam's second moment estimates to work properly.

The solution was LAMB\index{LAMB} (Layer-wise Adaptive Moments for Batch training)~\cite{you2019lamb}, which extends Adam with layer-wise normalization of updates, computing the ratio of parameter norm to update norm for each layer. This stabilizes training with very large batch sizes by preventing any single layer from dominating the update.

\textbf{Lesson:} Optimizers that work at small scale can fail at large scale. As you push the boundaries of batch size and model size, optimizer choice becomes critical. There's no universal ``best'' optimizer.
\end{warstory}

%------------------------------------------------------------------------------
\section{Learning Rate Scheduling}
\label{sec:lr-scheduling}

Rather than using a fixed learning rate throughout training, \textbf{learning rate schedules} adapt $\eta$ over time. This allows aggressive exploration early in training followed by refined convergence near the optimum.

\begin{principle}
A good learning rate schedule balances two competing goals:
\begin{enumerate}
    \item \textbf{Early training}: Large learning rates enable rapid progress and exploration
    \item \textbf{Late training}: Small learning rates enable fine-tuning and stable convergence
\end{enumerate}
\end{principle}

\subsection{Step Decay}

\begin{definition}{Step Decay Schedule}{step-decay-def}
Reduce the learning rate by a factor $\alpha$ every $k$ epochs:
\[
    \eta_t = \eta_0 \cdot \alpha^{\lfloor t / k \rfloor}
\]
Typical values: $\alpha = 0.1$ to $0.5$, $k = 10$ to $30$ epochs.
\end{definition}

Step decay is simple and effective, widely used for training CNNs on image classification tasks. The key is choosing when to drop the learning rate---typically when validation loss plateaus.

\subsection{Exponential Decay}

\begin{definition}{Exponential Decay Schedule}{exp-decay-def}
Decay the learning rate exponentially\index{exponential decay}:
\[
    \eta_t = \eta_0 \cdot e^{-\lambda t}
\]
where $\lambda > 0$ controls the decay rate.
\end{definition}

Exponential decay provides smooth, continuous reduction of the learning rate, avoiding the abrupt changes of step decay.

\subsection{Cosine Annealing}

\begin{definition}{Cosine Annealing Schedule}{cosine-def}
Follow a cosine curve from maximum to minimum learning rate\index{cosine annealing|textbf}:
\[
    \eta_t = \eta_{\min} + \frac{1}{2}(\eta_{\max} - \eta_{\min})\left(1 + \cos\left(\frac{\pi t}{T}\right)\right)
\]
where $T$ is the total number of training iterations.
\end{definition}

\begin{intuition}
Cosine annealing starts with a slow decrease, accelerates in the middle, and slows down again near the end. This profile matches the intuition that: (1) early training benefits from sustained high learning rates for exploration, (2) mid-training can tolerate faster decay as the model approaches good regions, and (3) late training needs careful fine-tuning with slowly decreasing rates.
\end{intuition}

Figure~\ref{fig:lr-schedules} compares the three schedules discussed in this section.

\begin{figure}[htbp]
\centering
\begin{tikzpicture}
\begin{axis}[
    width=0.85\textwidth,
    height=0.4\textwidth,
    xlabel={Epoch},
    ylabel={Learning Rate},
    xmin=0, xmax=100,
    ymin=0, ymax=0.12,
    grid=major,
    grid style={gray!20, thin},
    legend pos=north east,
    legend style={font=\small, fill=white, fill opacity=0.9, draw=gray!40, rounded corners=2pt},
    tick label style={font=\small},
    label style={font=\small},
    every axis plot/.append style={line width=1.2pt},
    samples=200
]
% Step decay
\addplot[primaryblue, domain=0:100, samples=100]
    {0.1 * 0.1^(floor(x/30))};
\addlegendentry{Step Decay}

% Exponential decay
\addplot[primaryred, domain=0:100]
    {0.1 * exp(-0.03*x)};
\addlegendentry{Exponential}

% Cosine annealing
\addplot[primarygreen, domain=0:100]
    {0.001 + 0.5*(0.1-0.001)*(1 + cos(deg(pi*x/100)))};
\addlegendentry{Cosine}

\end{axis}
\end{tikzpicture}
\caption{Comparison of learning rate schedules over 100 epochs. Step decay makes discrete jumps, exponential decay follows a smooth curve, and cosine annealing follows a smooth cosine curve that decays gradually.}
\label{fig:lr-schedules}
\end{figure}

\subsection{Cosine Annealing with Warm Restarts}

\begin{definition}{SGDR: Warm Restarts}{sgdr-def}
Cosine annealing with warm restarts (SGDR) periodically resets the learning rate to its initial value:
\[
    \eta_t = \eta_{\min} + \frac{1}{2}(\eta_{\max} - \eta_{\min})\left(1 + \cos\left(\frac{T_{\text{cur}}}{T_i}\pi\right)\right)
\]
where $T_{\text{cur}}$ is the number of epochs since the last restart, and $T_i$ is the period length. The period may increase after each restart: $T_{i+1} = T_i \cdot T_{\text{mult}}$.
\end{definition}

Warm restarts enable the optimizer to escape local minima by periodically ``reheating'' the learning rate. Each restart provides an opportunity to explore different regions of the loss landscape.

\subsection{Learning Rate Warmup}

\begin{definition}{Linear Warmup}{warmup-def}
Start with a small learning rate and linearly increase to the target value over $W$ warmup steps\index{learning rate!warmup}:
\[
    \eta_t = \begin{cases}
        \frac{t}{W} \cdot \eta_{\max} & \text{if } t < W \\
        \text{scheduled rate} & \text{otherwise}
    \end{cases}
\]
\end{definition}

\begin{keyconcept}
\textbf{Why Warmup Matters:}
\begin{itemize}
    \item At initialization, gradient estimates are unreliable and may have high variance
    \item Large initial learning rates can cause unstable updates, pushing parameters far from their initialized values
    \item Warmup allows the model to ``settle'' before taking larger steps
\end{itemize}
This is particularly important for Transformers and large batch training, where warmup is essential for stable convergence.
\end{keyconcept}

\begin{pythoncode}
import numpy as np

def warmup_cosine_schedule(epoch, total_epochs, warmup_epochs, initial_lr, min_lr=1e-6):
    """Warmup followed by cosine annealing schedule."""
    if epoch < warmup_epochs:
        return initial_lr * (epoch + 1) / warmup_epochs
    progress = (epoch - warmup_epochs) / (total_epochs - warmup_epochs)
    return min_lr + 0.5 * (initial_lr - min_lr) * (1 + np.cos(np.pi * progress))

def step_decay_schedule(epoch, initial_lr=0.1, drop_factor=0.1, drop_epochs=[30, 60, 90]):
    """Step decay schedule with configurable drop points."""
    return initial_lr * drop_factor ** sum(epoch >= e for e in drop_epochs)
\end{pythoncode}

\subsection{Cyclical Learning Rates}

\begin{definition}{Cyclical Learning Rate}{clr-def}
Oscillate the learning rate between a minimum and maximum value over a cycle of $2 \cdot \text{step\_size}$ iterations\index{cyclical learning rate|textbf}\index{learning rate!cyclical}:
\[
    \eta_t = \eta_{\min} + (\eta_{\max} - \eta_{\min}) \cdot \max(0, 1 - |x - 1|)
\]
where $x = t/\text{step\_size} - 2 \cdot \text{cycle} + 1$ and cycle $= \lfloor 1 + t/(2 \cdot \text{step\_size}) \rfloor$.
\end{definition}

Cyclical learning rates enable the model to periodically ``jump'' out of local minima. The triangular policy increases and decreases the learning rate linearly within each cycle. Variants include:
\begin{itemize}
    \item \textbf{Triangular}: Linear increase then decrease
    \item \textbf{Triangular2}: Same but cut the max learning rate in half each cycle
    \item \textbf{Exp\_range}: Exponential decay of the cycle amplitude
\end{itemize}

%------------------------------------------------------------------------------
\section{Optimizer Selection Guidelines}
\label{sec:optimizer-selection}

Choosing the right optimizer is problem-dependent. Here we provide practical guidelines based on empirical observations and theoretical understanding.

\begin{table}[htbp]
\centering
\caption{Optimizer selection guide by application domain}
\label{tab:optimizer-selection}
\begin{tabular}{@{}lllp{5cm}@{}}
\toprule
\textbf{Domain} & \textbf{Recommended} & \textbf{Alternative} & \textbf{Notes} \\
\midrule
CNNs (ImageNet) & SGD+Momentum & AdamW & Some studies suggest SGD may generalize better (Wilson et al., 2017), though results are task-dependent \\
Transformers & AdamW & Adam & Weight decay crucial \\
RNNs/LSTMs & Adam & RMSprop & Adaptive rates help \\
GANs & Adam & RMSprop & Lower $\beta_1$ (0.5) often used \\
Fine-tuning & AdamW & SGD & Small learning rate important \\
Quick prototyping & Adam & AdamW & Works well out-of-box \\
\bottomrule
\end{tabular}
\end{table}

\begin{principle}
\textbf{General Optimizer Selection Strategy:}
\begin{enumerate}
    \item Start with \textbf{AdamW} with default hyperparameters ($\eta = 0.001$, $\beta_1 = 0.9$, $\beta_2 = 0.999$, weight decay $= 0.01$)
    \item If training Transformers, use \textbf{warmup} (typically 10\% of total steps)
    \item For CNNs where generalization is critical, consider \textbf{SGD with momentum} ($\gamma = 0.9$)
    \item If Adam does not converge, try \textbf{RMSprop} or \textbf{reduce learning rate}
    \item For GANs, use Adam with $\beta_1 = 0.5$ to reduce momentum oscillations
\end{enumerate}
\end{principle}

\subsection{SGD vs. Adam: The Generalization Debate}

\begin{keyconcept}{SGD vs Adam: The Real Tradeoff}
You'll hear conflicting advice: ``Adam for everything'' vs. ``SGD generalizes better.'' Here's the nuanced truth:

\textbf{Adam wins on}:
\begin{itemize}
    \item Convergence speed (reaches good loss faster)
    \item Robustness to hyperparameters (less tuning needed)
    \item Transformers and attention architectures
    \item Quick prototyping and experimentation
\end{itemize}

\textbf{SGD (with momentum) wins on}:
\begin{itemize}
    \item Final generalization (slightly better test accuracy, when tuned)
    \item CNNs for image classification (ResNet, etc.)
    \item When you have compute budget to tune learning rate schedule
\end{itemize}

\textbf{The practical answer}: Start with Adam for fast iteration. Once you have a working model, try SGD+momentum with careful LR tuning for competitions or production where every 0.1\% matters.
\end{keyconcept}

\begin{warstory}
\textbf{Why AdamW Became Standard for Transformer Training}

The computer vision community long preferred SGD with momentum for final accuracy. But when practitioners began training large Transformer models, a different pattern emerged.

Empirical results suggest that AdamW often outperforms SGD on Transformer architectures. One hypothesis for why: Transformers have heterogeneous parameter types---embedding matrices, attention weights, and feed-forward layers---each with different gradient statistics. Adam's per-parameter learning rates may better adapt to this heterogeneity, while SGD's single learning rate cannot optimize all components equally.

The shift to \textbf{AdamW} (decoupled weight decay) rather than Adam with L2 regularization followed Loshchilov and Hutter's 2019 paper showing that standard Adam's L2 regularization gets scaled by the adaptive learning rate, effectively \emph{reducing} regularization for parameters with large gradients. AdamW applies weight decay directly, providing more consistent regularization.

Today, AdamW with warmup and cosine decay is a common recipe for large language model training, though specific configurations vary across organizations and models.

\textbf{Lesson:} Optimizer choice often depends on architecture. SGD tends to work well for CNNs; AdamW is popular for Transformers. But always validate on your specific task.
\end{warstory}

There is ongoing debate about whether adaptive optimizers like Adam find solutions that generalize as well as SGD. Several observations inform this discussion:

\begin{itemize}
    \item \textbf{Flat minima hypothesis}\index{flat minima}: SGD with momentum tends to find wider minima that generalize better. Adam's adaptive rates may lead to sharper minima\index{sharp minima}.
    \item \textbf{AdamW improvements}: Decoupling weight decay significantly improves Adam's generalization, narrowing the gap with SGD.
    \item \textbf{Task dependence}: NLP tasks with Transformers tend to favor Adam/AdamW, while computer vision with CNNs often favors SGD, though results vary.
    \item \textbf{Learning rate tuning}: Properly tuned SGD can match or exceed Adam's performance, but requires more careful hyperparameter search.
\end{itemize}

\begin{perspective}
The choice between SGD and Adam is not about which is ``better'' but about trade-offs. Adam tends to converge faster and is generally less sensitive to hyperparameters, making it well-suited for prototyping and domains where fast iteration matters. SGD may achieve marginally better final performance with careful tuning, which can be important for competitive benchmarks. In practice, both optimizers can produce models that perform well when properly configured.
\end{perspective}

\begin{checkpoint}
When would you choose SGD over Adam? Give two scenarios with reasoning.
\end{checkpoint}

\begin{table}[htbp]
\centering
\caption{Optimizer Selection Guide: When to Use What}
\label{tab:optimizer-selection-guide}
\begin{tabular}{@{}llll@{}}
\toprule
\textbf{Situation} & \textbf{Best Choice} & \textbf{Avoid} & \textbf{Reason} \\
\midrule
Computer Vision (CNNs) & SGD + Momentum & Adam & Better generalization \\
NLP / Transformers & AdamW & SGD & Handles sparse gradients \\
Small dataset (<10k) & SGD & Adam & Less overfitting risk \\
Fine-tuning pretrained & AdamW, low LR & SGD & Stability matters \\
Reinforcement Learning & Adam & SGD & Non-stationary objectives \\
GANs & Adam ($\beta_1=0.5$) & SGD & Training stability \\
Very deep networks & Adam/LAMB & Vanilla SGD & Adaptive LR per layer \\
Convex problems & SGD & Adam & Provable convergence \\
\bottomrule
\end{tabular}
\end{table}

%------------------------------------------------------------------------------
\subsection{Visual Comparison: Optimizers on an Ill-Conditioned Problem}

To make these differences concrete, let's trace each optimizer on the same ill-conditioned quadratic:
\begin{equation}
L(\theta_1, \theta_2) = 0.5\theta_1^2 + 25\theta_2^2
\end{equation}
The condition number $\kappa = 50$ creates elongated contours---$\theta_2$ has much steeper gradients than $\theta_1$. All optimizers start at $\theta_0 = (4, 1)$ with learning rate $\eta = 0.02$, as shown in Figure~\ref{fig:optimizer-comparison}.

\begin{figure}[htbp]
\centering
\begin{tikzpicture}[scale=0.9]
% Draw contours of ill-conditioned quadratic
\begin{axis}[
    view={0}{90},
    width=11cm,
    height=8cm,
    xlabel={$\theta_1$},
    ylabel={$\theta_2$},
    xmin=-1, xmax=5,
    ymin=-1.2, ymax=1.5,
    axis lines=middle,
    xtick={0,2,4},
    ytick={-1,0,1},
    tick label style={font=\small},
    label style={font=\small},
    clip=false
]

% Contour ellipses (0.5*x^2 + 25*y^2 = c for various c)
% c = 0.5: x^2 + 50y^2 = 1, so a=1, b=0.141
% c = 2: x^2 + 50y^2 = 4, so a=2, b=0.283
% c = 4.5: x^2 + 50y^2 = 9, so a=3, b=0.424
% c = 8: x^2 + 50y^2 = 16, so a=4, b=0.566
% c = 12.5: x^2 + 50y^2 = 25, so a=5, b=0.707

\draw[gray!20, thin] (0,0) ellipse (1 and 0.141);
\draw[gray!20, thin] (0,0) ellipse (2 and 0.283);
\draw[gray!20, thin] (0,0) ellipse (3 and 0.424);
\draw[gray!20, thin] (0,0) ellipse (4 and 0.566);

% Note: Trajectories are illustrative approximations for pedagogical clarity
% SGD trajectory (oscillates badly due to high gradient in y direction)
\addplot[primaryblue, line width=1.5pt, ->, mark=*, mark size=2pt] coordinates {
    (4.0, 1.0) (3.92, 0.0) (3.84, 0.0) (3.77, 0.0) (3.69, 0.0)
};
\node[primaryblue, right, font=\footnotesize\bfseries] at (axis cs:3.7, 0.15) {SGD};

% Momentum trajectory (overshoots then oscillates back)
\addplot[primaryred, line width=1.5pt, ->, mark=*, mark size=2pt] coordinates {
    (4.0, 1.0) (3.92, -0.9) (3.74, 0.72) (3.48, -0.50) (3.17, 0.28)
    (2.84, -0.12) (2.52, 0.04)
};
\node[primaryred, below, font=\footnotesize\bfseries] at (axis cs:2.5, -0.95) {Momentum};

% RMSprop trajectory (adapts per-coordinate, more direct)
\addplot[accentorange, line width=1.5pt, ->, mark=*, mark size=2pt] coordinates {
    (4.0, 1.0) (3.55, 0.55) (3.15, 0.30) (2.80, 0.16) (2.48, 0.09)
    (2.20, 0.05) (1.95, 0.02)
};
\node[accentorange, above right, font=\footnotesize\bfseries] at (axis cs:2.5, 0.25) {RMSprop};

% Adam trajectory (combines benefits, smooth convergence)
\addplot[primarygreen, line width=1.5pt, ->, mark=*, mark size=2pt] coordinates {
    (4.0, 1.0) (3.50, 0.50) (3.05, 0.24) (2.65, 0.11) (2.30, 0.05)
    (1.98, 0.02) (1.70, 0.01) (1.45, 0.005)
};
\node[primarygreen, below left, font=\footnotesize\bfseries] at (axis cs:1.5, 0.08) {Adam};

% Mark minimum
\node[star, star points=5, fill=black, inner sep=1.5pt] at (axis cs:0, 0) {};
\node[below right] at (axis cs:0.1, -0.1) {$\theta^*$};

% Mark start
\node[circle, fill=black, inner sep=2pt] at (axis cs:4, 1) {};
\node[above right] at (axis cs:4, 1) {Start};

\end{axis}
\end{tikzpicture}
\caption{Optimizer trajectories on an ill-conditioned quadratic with condition number $\kappa=50$. \textbf{SGD} (blue) makes slow progress in $\theta_1$ while rapidly collapsing $\theta_2$ to zero. \textbf{Momentum} (red) accelerates but oscillates in $\theta_2$. \textbf{RMSprop} (orange) adapts step sizes per-coordinate, reducing oscillation. \textbf{Adam} (green) combines both benefits for smooth, rapid convergence. \textbf{Note:} This is a schematic illustration; actual trajectories depend on specific hyperparameter choices.}
\label{fig:optimizer-comparison}
\end{figure}

\begin{example}{Optimizer Behavior Analysis}{optimizer-trace}
Let's trace the first few steps analytically. The gradient is $\nabla L = (\theta_1, 50\theta_2)$.

\textbf{Vanilla SGD} ($\eta = 0.02$): At $\theta_0 = (4, 1)$, the gradient is $(4, 50)$. The update:
\[
\theta_1 = (4 - 0.02 \cdot 4,\; 1 - 0.02 \cdot 50) = (3.92,\; 0)
\]
The step in $\theta_2$ is exactly $-1$, landing precisely at zero. But the step in $\theta_1$ is tiny (0.08). SGD makes fast progress in the steep direction but crawls in the shallow one.

\textbf{SGD with Momentum} ($\gamma = 0.9$): Momentum accumulates velocity. After several steps, the oscillations in $\theta_2$ build up velocity that overshoots, while $\theta_1$ gradually accelerates. The trajectory zigzags.

\textbf{RMSprop} ($\beta = 0.9$): Maintains running averages $v_1$ and $v_2$ of squared gradients. Since $\theta_2$ gradients are $\sim 50\times$ larger, $\sqrt{v_2} \gg \sqrt{v_1}$. The effective step sizes equalize:
\[
\frac{\eta}{\sqrt{v_1 + \epsilon}} \approx \frac{\eta}{\sqrt{v_2 + \epsilon}}
\]
This ``condition number repair'' makes the trajectory more direct toward the minimum.

\textbf{Adam}: Combines RMSprop's adaptive rates with momentum's acceleration. The bias-corrected updates provide both equalized step sizes and smooth accumulation of direction, typically yielding a more efficient path on ill-conditioned problems.

\textit{Note: This analysis illustrates characteristic optimizer behaviors on a specific quadratic; actual performance varies with problem structure and hyperparameters.}
\end{example}

\begin{keyconcept}
\textbf{Why Adaptive Methods Excel on Ill-Conditioned Problems.}
The fundamental insight is \textbf{per-parameter learning rates}\index{per-parameter learning rate|textbf}. When features have vastly different scales (common in neural networks with embeddings, batch norm parameters, and output layers), a single global learning rate cannot be optimal for all parameters simultaneously. Adaptive methods automatically tune effective learning rates based on gradient history, achieving what would otherwise require careful per-layer learning rate schedules.
\end{keyconcept}

%------------------------------------------------------------------------------
\section{Learning Rate Selection in Practice}\label{sec:lr-finder}

Theory tells us the learning rate should satisfy $\eta \leq \frac{1}{L}$ where $L$ is the smoothness constant. But how do you find $L$ for a neural network? You don't---it depends on the weights, data, and architecture in complex ways.

Instead, practitioners use empirical methods.

\begin{principle}{The Learning Rate Finder}
The learning rate finder (popularized by fast.ai) sweeps through learning rates exponentially and plots the loss:
\begin{enumerate}
    \item Start with a tiny LR ($10^{-8}$) and train for one batch
    \item Multiply LR by a factor (e.g., 1.1) and train another batch
    \item Repeat until loss explodes or you reach a large LR ($10$)
    \item Plot loss vs. learning rate (log scale)
    \item Choose the LR where loss decreases fastest (steepest negative slope), typically 1/10th of the LR where loss starts increasing
\end{enumerate}
\end{principle}

\begin{lstlisting}[language=Python, caption={Learning rate finder pseudocode}]
def find_lr(model, loader, start_lr=1e-8, end_lr=10, num_steps=100):
    lrs, losses = [], []
    lr = start_lr
    factor = (end_lr / start_lr) ** (1 / num_steps)

    for step, (x, y) in enumerate(loader):
        if step >= num_steps:
            break
        set_lr(optimizer, lr)
        loss = train_step(model, x, y)
        lrs.append(lr)
        losses.append(loss)
        if loss > 4 * min(losses):  # Loss exploding
            break
        lr *= factor

    return lrs, losses  # Plot these to find optimal LR
\end{lstlisting}

%------------------------------------------------------------------------------
\section{Practical Hyperparameter Tuning}
\label{sec:hyperparameter-tuning}

\subsection{Learning Rate Finding}

\begin{definition}{Learning Rate Range Test}{lr-range-def}
Start with a very small learning rate and exponentially increase it over a mini-training run. Plot loss vs. learning rate and select the rate where loss decreases most steeply (typically just before the loss starts to increase).
\end{definition}

\begin{algorithm}[htbp]
\caption{Learning Rate Range Test}
\label{alg:lr-range}
\begin{algorithmic}[1]
\Require Model, data loader, $\eta_{\min} = 10^{-8}$, $\eta_{\max} = 1$, num\_iters
\Ensure Plot of loss vs. learning rate
\State $\eta_t \gets \eta_{\min}$
\State $\text{mult} \gets (\eta_{\max}/\eta_{\min})^{1/\text{num\_iters}}$
\For{$t = 1$ to num\_iters}
    \State Train one batch with learning rate $\eta_t$
    \State Record loss
    \State $\eta_t \gets \eta_t \cdot \text{mult}$
    \If{loss diverges (e.g., $> 4 \times$ initial loss)}
        \State \textbf{break}
    \EndIf
\EndFor
\State Plot smoothed loss vs. log learning rate
\end{algorithmic}
\end{algorithm}

\begin{pythoncode}
def lr_range_test(model, train_loader, min_lr=1e-8, max_lr=1, num_iter=100):
    """Perform learning rate range test."""
    mult = (max_lr / min_lr) ** (1 / num_iter)
    lr, lrs, losses, best_loss = min_lr, [], [], float('inf')

    for i, (inputs, targets) in enumerate(train_loader):
        if i >= num_iter:
            break
        for pg in optimizer.param_groups:
            pg['lr'] = lr
        optimizer.zero_grad()
        loss = criterion(model(inputs), targets)
        loss.backward()
        optimizer.step()
        lrs.append(lr)
        losses.append(loss.item())
        if loss.item() > 4 * best_loss:
            break
        best_loss = min(best_loss, loss.item())
        lr *= mult
    return lrs, losses
\end{pythoncode}

\subsection{Batch Size and Learning Rate Scaling}

When increasing batch size, the learning rate should often be scaled proportionally to maintain similar training dynamics.

\begin{principle}[title=Linear Scaling Rule]
\textbf{(Empirical Heuristic)} When multiplying the batch size by $k$, multiply the learning rate by $k$:
\[
    \eta_{\text{new}} = k \cdot \eta_{\text{base}}
\]
This approximately preserves the expected weight update magnitude per epoch. Note: This is an empirical guideline (Goyal et al., 2017), not a theorem---it breaks down for very large batch sizes and requires warmup to work reliably.
\end{principle}

\begin{important}
The linear scaling rule has limitations:
\begin{itemize}
    \item Requires warmup when scaling to very large batch sizes
    \item Does not hold for batch sizes beyond a certain threshold
    \item Square root scaling ($\eta_{\text{new}} = \sqrt{k} \cdot \eta_{\text{base}}$) may be more stable
\end{itemize}
\end{important}

\subsection{Hyperparameter Search Strategies}

\begin{definition}{Grid Search}{grid-search-def}
Exhaustively evaluate all combinations of hyperparameters from predefined discrete sets\index{grid search}:
\[
    \text{Grid} = \{\eta_1, \ldots, \eta_m\} \times \{\gamma_1, \ldots, \gamma_n\} \times \cdots
\]
\end{definition}

\begin{definition}{Random Search}{random-search-def}
Sample hyperparameters randomly from specified distributions (uniform, log-uniform, etc.). More efficient than grid search for high-dimensional search spaces\index{random search}\index{hyperparameter search}.
\end{definition}

\begin{theorem}{Random Search Efficiency}{random-efficiency}
For hyperparameter optimization where only a few hyperparameters significantly affect performance, random search finds good configurations with fewer evaluations than grid search. Specifically, if $k$ out of $n$ hyperparameters are important, random search with $m$ trials has probability $1 - (1 - 3^{-k})^m$ of finding a configuration in the top $3^{-k}$ fraction.
\end{theorem}

\begin{pythoncode}
from scipy.stats import loguniform

def random_hyperparameter_search(train_fn, n_trials=20):
    """Random hyperparameter search for optimizer settings."""
    best_loss, best_params = float('inf'), None
    for trial in range(n_trials):
        params = {
            'learning_rate': loguniform.rvs(1e-5, 1e-2),
            'weight_decay': loguniform.rvs(1e-6, 1e-2),
            'beta1': np.random.uniform(0.85, 0.95),
            'beta2': np.random.uniform(0.99, 0.9999),
        }
        val_loss = train_fn(**params)
        if val_loss < best_loss:
            best_loss, best_params = val_loss, params
    return best_params, best_loss
\end{pythoncode}

%------------------------------------------------------------------------------
\section{Common Pitfalls and Solutions}
\label{sec:pitfalls}

\subsection{Gradient Explosion and Clipping}

Large gradients can destabilize training, particularly in RNNs and deep networks. \textbf{Gradient clipping}\index{gradient scaling|see{gradient clipping}} constrains gradient magnitudes.

\begin{definition}{Gradient Clipping by Norm}{grad-clip-def}
Scale gradients so their global norm does not exceed a threshold $c$\index{gradient clipping|textbf}:
\[
    \vect{g} \gets \begin{cases}
        \vect{g} \cdot \frac{c}{\norm{\vect{g}}} & \text{if } \norm{\vect{g}} > c \\
        \vect{g} & \text{otherwise}
    \end{cases}
\]
Typical values: $c = 1.0$ to $5.0$.
\end{definition}

\begin{pythoncode}
@torch.no_grad()
def clip_grad_norm(parameters, max_norm):
    """Clip gradients by global norm."""
    params = [p for p in parameters if p.grad is not None]
    total_norm = torch.stack([p.grad.norm(2) for p in params]).norm(2)
    clip_coef = max_norm / (total_norm + 1e-6)
    if clip_coef < 1:
        for p in params:
            p.grad.mul_(clip_coef)
    return total_norm.item()
\end{pythoncode}

\subsection{Learning Rate Warmup for Stability}

Training instabilities at the start can derail the entire training run. Warmup provides a safety net.

\begin{important}
\textbf{When Warmup is Essential:}
\begin{itemize}
    \item Training Transformers (especially large models)
    \item Using large batch sizes
    \item After initialization schemes that may produce large gradients
    \item When combining with aggressive learning rate schedules
\end{itemize}
A common heuristic: warmup for 5--10\% of total training steps.
\end{important}

\subsection{Numerical Precision Issues}

Adam's bias-corrected second moment $\hat{v}_t$ can become very small for some parameters, causing numerical issues when dividing.

\begin{itemize}
    \item \textbf{Solution 1}: Use $\epsilon$ in denominator (already part of Adam)
    \item \textbf{Solution 2}: Use mixed precision training with loss scaling
    \item \textbf{Solution 3}: Increase $\epsilon$ for particularly problematic parameters
\end{itemize}

\begin{whatbreaks}
\textbf{Exercise 1:} Remove bias correction from Adam (use $m_t$ and $v_t$ directly instead of $\hat{m}_t$ and $\hat{v}_t$). Compare the first 5 parameter updates to corrected Adam. What's different?

\textbf{Exercise 2:} Set $\beta_2 = 0$ in Adam (disable second moment). What algorithm do you get? Does it still converge? Test on a simple quadratic.
\end{whatbreaks}

%------------------------------------------------------------------------------
\section{Diagnosing Adam Failure Modes}
\label{sec:adam-failure-modes}

Adam is remarkably robust, but it can fail in specific situations. Understanding these failure modes helps diagnose training problems and choose appropriate alternatives.

\subsection{The ``Adam Doesn't Generalize'' Problem}

\begin{important}
\textbf{Symptom: Adam achieves low training loss but test performance is worse than SGD}

\textbf{Why it happens:} Adam's adaptive learning rates can converge to sharp minima that generalize poorly. The per-parameter scaling can also interfere with implicit regularization from weight decay.

\textbf{Diagnosis:}
\begin{enumerate}
    \item Compare training and test loss curves between Adam and SGD
    \item If Adam's training loss is lower but test loss is higher, this is the generalization gap
    \item Compute the Hessian's maximum eigenvalue at convergence---Adam often finds sharper minima
\end{enumerate}

\textbf{Fixes (in order of effectiveness):}
\begin{enumerate}
    \item \textbf{Switch to AdamW}: Decoupled weight decay often closes the gap with SGD
    \item \textbf{Increase weight decay}: Try $\lambda = 0.01$ to $0.1$ (higher than SGD defaults)
    \item \textbf{Reduce $\beta_2$}: Values like $0.99$ instead of $0.999$ reduce ``memory'' of past gradients
    \item \textbf{Add explicit regularization}: Dropout, label smoothing, data augmentation
    \item \textbf{Use SGD for final fine-tuning}: Train with Adam, then switch to SGD for last 10--20\% of training
\end{enumerate}
\end{important}

\begin{keyconcept}
\textbf{Why AdamW Helps Generalization:}

Standard Adam with L2 regularization applies the penalty \textit{before} adaptive scaling:
\[
    \vect{\theta}_{t+1} = \vect{\theta}_t - \frac{\eta}{\sqrt{\hat{\vect{v}}_t} + \epsilon} (\hat{\vect{m}}_t + \lambda \vect{\theta}_t)
\]

The weight decay term $\lambda \vect{\theta}_t$ is scaled by $1/\sqrt{\hat{\vect{v}}_t}$, which varies per parameter. Parameters with small gradients receive \textit{stronger} regularization than intended.

AdamW applies weight decay \textit{after} adaptive scaling:
\[
    \vect{\theta}_{t+1} = \vect{\theta}_t - \eta \left( \frac{\hat{\vect{m}}_t}{\sqrt{\hat{\vect{v}}_t} + \epsilon} + \lambda \vect{\theta}_t \right)
\]

Now all parameters receive uniform regularization strength $\eta \lambda$, matching the behavior of SGD with weight decay.
\end{keyconcept}

\subsection{Adaptive Methods and Sparse Gradients}

\begin{important}
\textbf{Symptom: Adam performs poorly on embeddings or sparse feature problems}

\textbf{Why it happens:} For sparse features (e.g., word embeddings where most words appear rarely), Adam's second moment $\vect{v}_t$ accumulates slowly for rare features. With $\beta_2 = 0.999$, it takes thousands of updates to get reliable estimates. During this time, rare features receive inappropriately large updates.

\textbf{Diagnosis:}
\begin{enumerate}
    \item Check if your problem has sparse gradients (embeddings, one-hot inputs, sparse attention)
    \item Monitor update magnitudes for frequent vs. rare features---large disparities indicate the problem
    \item Compare Adam to SGD on the sparse components specifically
\end{enumerate}

\textbf{Fixes:}
\begin{enumerate}
    \item \textbf{Use lower $\beta_2$}: Try $0.99$ or $0.9$ for faster adaptation to sparse features
    \item \textbf{Use AMSGrad}: Variant that maintains maximum of past $\vect{v}_t$ values, preventing rare-feature explosion
    \item \textbf{Separate optimizers}: Use Adam for dense layers, SGD for embeddings
    \item \textbf{Use AdaFactor}: Designed specifically for large embedding tables
\end{enumerate}
\end{important}

\begin{commonmistakes}
\begin{itemize}
    \item \textbf{Mistake:} Using Adam's default $\beta_2 = 0.999$ for all problems

    \textbf{Why it's wrong:} The default assumes gradients are relatively stationary. For non-stationary problems (RL, GANs, curriculum learning) or sparse problems (NLP embeddings), $\beta_2 = 0.999$ means Adam ``remembers'' irrelevant old gradients for too long.

    \textbf{Correct approach:}
    \begin{itemize}
        \item Stationary, dense problems: $\beta_2 = 0.999$ (default)
        \item Sparse features: $\beta_2 = 0.99$
        \item Non-stationary (RL, GANs): $\beta_2 = 0.99$ or lower
        \item Very long training: Consider $\beta_2 = 0.9999$ to reduce noise
    \end{itemize}
\end{itemize}
\end{commonmistakes}

\subsection{Epsilon ($\epsilon$) Too Small or Too Large}

\begin{important}
\textbf{Epsilon ($\epsilon$) Tuning:}

\textbf{NaN or Inf values appear:} $\epsilon$ is too small. For parameters with very small gradients, $\hat{\vect{v}}_t$ can underflow, causing astronomically large updates. \emph{Fix:} Increase $\epsilon$ from $10^{-8}$ to $10^{-6}$ or $10^{-4}$. For mixed-precision training, $\epsilon = 10^{-4}$ is often necessary.

\textbf{Training is stable but slow for some parameters:} $\epsilon$ is too large. When $\sqrt{\hat{\vect{v}}_t} \ll \epsilon$, the effective learning rate becomes $\eta / \epsilon$ regardless of gradient history, and Adam degrades to near-SGD behavior. \emph{Fix:} Reduce $\epsilon$ to $10^{-8}$ (default). If you increased $\epsilon$ for stability, consider gradient clipping instead.
\end{important}

\begin{keyconcept}
\textbf{Epsilon Selection Guide:}

\begin{center}
\begin{tabular}{ll}
\toprule
\textbf{Scenario} & \textbf{Recommended $\epsilon$} \\
\midrule
Default / most problems & $10^{-8}$ \\
Mixed precision (FP16) & $10^{-4}$ to $10^{-3}$ \\
Very deep networks (100+ layers) & $10^{-6}$ \\
Embedding layers with rare tokens & $10^{-6}$ \\
If seeing NaN and clipping doesn't help & Increase by $10\times$ \\
\bottomrule
\end{tabular}
\end{center}

\textbf{Rule of thumb:} Start with $\epsilon = 10^{-8}$. If you see NaN, first try gradient clipping. If that fails, increase $\epsilon$ to $10^{-6}$, then $10^{-4}$.
\end{keyconcept}

\subsection{Adam Diagnostic Checklist}

\begin{principle}
\textbf{When Adam Training Fails or Underperforms:}

\begin{enumerate}
    \item \textbf{Loss is NaN or Inf?}
    \begin{itemize}
        \item Add gradient clipping at $\norm{\vect{g}} \leq 1.0$
        \item Increase $\epsilon$ to $10^{-6}$ or $10^{-4}$
        \item Check for numerical issues in loss computation (log of zero, etc.)
    \end{itemize}

    \item \textbf{Training loss decreases but test loss doesn't?}
    \begin{itemize}
        \item Switch to AdamW with weight decay $\lambda = 0.01$
        \item Reduce $\beta_2$ to $0.99$
        \item Consider switching to SGD+Momentum for final training phase
    \end{itemize}

    \item \textbf{Some parameters train well, others don't?}
    \begin{itemize}
        \item Check for sparse gradient patterns
        \item Use per-layer learning rates (smaller for pretrained, larger for new)
        \item Try AMSGrad or AdaFactor for embedding-heavy models
    \end{itemize}

    \item \textbf{Training is unstable early, then stabilizes?}
    \begin{itemize}
        \item Add learning rate warmup (5--10\% of training)
        \item The bias correction is struggling---warmup gives it time
    \end{itemize}

    \item \textbf{Training stalls after initial progress?}
    \begin{itemize}
        \item Second moment may have ``memorized'' early noise
        \item Reduce $\beta_2$ to $0.99$ for faster forgetting
        \item Try a warmup restart: reset optimizer state and warmup again
    \end{itemize}
\end{enumerate}
\end{principle}

\begin{warstory}
\textbf{The Embedding That Broke Adam}

A team training a recommendation system noticed that common items (viewed millions of times) learned well, but rare items (viewed dozens of times) had erratic embeddings. The model would occasionally recommend bizarre items to users.

The diagnosis: with $\beta_2 = 0.999$, Adam's second moment for rare item embeddings was dominated by the initialization phase. The few updates these embeddings received were scaled by near-zero $\sqrt{\hat{\vect{v}}_t}$, producing enormous effective learning rates. Rare items' embeddings were essentially random walks.

The fix was threefold: (1) Reduce $\beta_2$ to $0.99$ for the embedding layer, (2) Use gradient clipping specifically on embedding gradients, and (3) Initialize $\vect{v}_0$ to a small positive value (0.01) instead of zero for embeddings.

\textbf{Lesson:} Adam's defaults assume all parameters see gradients at similar rates. For embedding tables with power-law popularity distributions, this assumption fails catastrophically.
\end{warstory}

%------------------------------------------------------------------------------
\section{Exercises}
\label{sec:adv-opt-exercises}

\begin{exercise}{Adaptive Methods vs. Momentum}{ex-adaptive-analysis}
Consider minimizing $f(x, y) = x^2 + 100y^2$ starting from $(10, 1)$. This ill-conditioned problem has different curvatures along the two axes.
\begin{enumerate}[label=(\alph*)]
    \item Run vanilla gradient descent with $\eta = 0.01$ for 500 iterations. Plot the trajectory.
    \item Run SGD with momentum ($\gamma = 0.9$, see Chapter~\ref{ch:gradient-descent}). How does the trajectory change?
    \item Run Adam with default hyperparameters. Compare to the momentum trajectory.
    \item Explain why Adam handles this ill-conditioned problem better than momentum alone.
\end{enumerate}
\end{exercise}

\begin{exercise}{Adam Bias Correction}{ex-adam-bias}
\begin{enumerate}[label=(\alph*)]
    \item Derive the expected value $\E[\vect{m}_t]$ assuming $\E[\vect{g}_i] = \vect{\mu}$ for all $i$.
    \item Show that the bias-corrected estimate $\hat{\vect{m}}_t = \vect{m}_t / (1 - \beta_1^t)$ is unbiased.
    \item Plot the bias correction factor $1/(1 - \beta_1^t)$ for $\beta_1 = 0.9$ over 100 iterations.
    \item At what iteration does the correction become negligible (say, within 1\% of 1)?
\end{enumerate}
\end{exercise}

\begin{exercise}{Learning Rate Schedule Comparison}{ex-lr-schedule}
Train a simple CNN on CIFAR-10 with three different learning rate schedules:
\begin{enumerate}[label=(\alph*)]
    \item Constant learning rate $\eta = 0.1$
    \item Step decay: multiply by 0.1 at epochs 30 and 60
    \item Cosine annealing with warmup (5 epoch warmup)
\end{enumerate}
Compare training curves and final test accuracy. Which schedule performs best and why?
\end{exercise}

\begin{exercise}{AdamW vs. Adam + L2}{ex-adamw}
\begin{enumerate}[label=(\alph*)]
    \item Mathematically show that L2 regularization in Adam is not equivalent to weight decay.
    \item Design an experiment to demonstrate the practical difference.
    \item Under what conditions would you expect the difference to be most pronounced?
\end{enumerate}
\end{exercise}

\begin{exercise}{Optimizer Implementation}{ex-impl}
Implement the following optimizers from scratch in NumPy:
\begin{enumerate}[label=(\alph*)]
    \item RMSprop
    \item Adam
    \item AdamW
\end{enumerate}
Test your implementations on the Rosenbrock function $f(x, y) = (1-x)^2 + 100(y-x^2)^2$ and compare convergence trajectories. For comparison with momentum-based methods, see the exercises in Chapter~\ref{ch:gradient-descent}.
\end{exercise}

%------------------------------------------------------------------------------
\section{Summary}
\label{sec:adv-opt-summary}

Building on the momentum methods covered in Chapter~\ref{ch:gradient-descent}, this chapter explored adaptive learning rate methods that power modern deep learning:

\begin{itemize}
    \item \textbf{Adaptive learning rate methods} automatically scale the learning rate for each parameter based on gradient history. AdaGrad pioneered this approach, RMSprop addressed its diminishing learning rate problem, and Adam combined adaptivity with momentum.

    \item \textbf{AdamW} decouples weight decay from the gradient update, improving generalization compared to standard L2 regularization in Adam.

    \item \textbf{Learning rate schedules} adapt the learning rate over training. Warmup stabilizes early training, while decay schedules (step, exponential, cosine) enable fine-tuning near convergence.

    \item \textbf{Optimizer selection} depends on the problem domain: SGD with momentum (Chapter~\ref{ch:gradient-descent}) often generalizes better for CNNs, while AdamW is preferred for Transformers.

    \item \textbf{Hyperparameter tuning} requires systematic approaches. Learning rate range tests identify good initial rates, and random search efficiently explores hyperparameter spaces.
\end{itemize}

The next chapter addresses the other half of successful training: \textbf{regularization}---the techniques that prevent overfitting and ensure our optimized models generalize to new data.

\section*{Interview Questions}
\begin{interviewquestions}

\textbf{Question 1: Explain Adam. Why has it become so popular?}

\textit{Expected follow-ups:}
\begin{itemize}
    \item What are the first and second moments tracking?
    \item Why is bias correction needed?
    \item What are the default hyperparameter values and why?
\end{itemize}

\textit{Red flags to avoid:}
\begin{itemize}
    \item Not mentioning bias correction---it's crucial for understanding Adam
    \item Saying Adam is ``always better than SGD'' (it's not for CNNs on ImageNet)
    \item Confusing the roles of $\beta_1$ and $\beta_2$
\end{itemize}

\textit{Strong answer key points:}
\begin{itemize}
    \item \textbf{Adam = Momentum + RMSprop + Bias Correction}: First moment $m_t$ tracks gradient direction (momentum), second moment $v_t$ tracks gradient magnitude (adaptive LR)
    \item \textbf{Update}: $\theta_{t+1} = \theta_t - \eta \cdot \hat{m}_t / (\sqrt{\hat{v}_t} + \epsilon)$ where $\hat{m}_t, \hat{v}_t$ are bias-corrected
    \item \textbf{Bias correction}: Since $m_0 = v_0 = 0$, early estimates are biased toward zero. Divide by $(1 - \beta^t)$ to correct
    \item \textbf{Why popular}: Works well out-of-box with minimal tuning; default $\eta=0.001$, $\beta_1=0.9$, $\beta_2=0.999$ work for most problems
    \item \textbf{Implicit trust region}: The $\sqrt{v_t}$ denominator effectively bounds step size regardless of gradient scale
\end{itemize}

\vspace{0.5em}
\textbf{Question 2: Why might Adam not generalize as well as SGD?}

\textit{Expected follow-ups:}
\begin{itemize}
    \item What evidence supports this claim?
    \item Does AdamW fix this issue?
    \item When would you still choose Adam over SGD?
\end{itemize}

\textit{Red flags to avoid:}
\begin{itemize}
    \item Saying ``Adam always generalizes worse''---it depends heavily on task and architecture
    \item Not distinguishing between Adam and AdamW
    \item Not acknowledging that this is an active area of research with nuanced findings
\end{itemize}

\textit{Strong answer key points:}
\begin{itemize}
    \item \textbf{Sharp minima hypothesis}: Adam may converge to sharper minima than SGD; sharper minima generalize worse to test data
    \item \textbf{L2 vs weight decay}: In Adam (not AdamW), L2 regularization is scaled by adaptive learning rates, weakening regularization for parameters with large gradients
    \item \textbf{Empirical evidence}: On ImageNet with ResNets, SGD+momentum typically achieves 0.5-1\% better top-1 accuracy
    \item \textbf{Counter-examples}: For Transformers and NLP tasks, Adam/AdamW consistently outperforms SGD
    \item \textbf{When to use Adam}: Rapid prototyping, Transformers, RNNs, any setting where tuning time is valuable
    \item \textbf{AdamW partially fixes this}: Decoupled weight decay maintains consistent regularization across parameters
\end{itemize}

\vspace{0.5em}
\textbf{Question 3: What's the difference between Adam and AdamW? When does it matter?}

\textit{Expected follow-ups:}
\begin{itemize}
    \item Can you write out both update rules?
    \item Why does this distinction matter for Transformers specifically?
    \item Is one always better than the other?
\end{itemize}

\textit{Red flags to avoid:}
\begin{itemize}
    \item Not knowing the difference (very common in interviews!)
    \item Saying they're ``basically the same''
    \item Not understanding that L2 regularization and weight decay are only equivalent for vanilla SGD
\end{itemize}

\textit{Strong answer key points:}
\begin{itemize}
    \item \textbf{L2 in Adam}: Adds $\lambda\theta$ to gradient \textit{before} adaptive scaling: $g_t = \nabla L + \lambda\theta$, then apply Adam update
    \item \textbf{Weight decay in AdamW}: Applies decay \textit{after} adaptive scaling: $\theta_{t+1} = \theta_t - \eta(\hat{m}_t/(\sqrt{\hat{v}_t}+\epsilon) + \lambda\theta_t)$
    \item \textbf{Why it matters}: With L2 in Adam, regularization strength varies per-parameter based on $\sqrt{v_t}$. Parameters with large gradients get weaker regularization
    \item \textbf{For Transformers}: Attention weights and embedding layers have different gradient magnitudes; AdamW ensures consistent regularization
    \item \textbf{Rule of thumb}: Always use AdamW for modern architectures; Adam with L2 is a legacy pattern
\end{itemize}

\vspace{0.5em}
\textbf{Question 4: When would you choose SGD over Adam?}

\textit{Expected follow-ups:}
\begin{itemize}
    \item What about for Transformers?
    \item How much tuning does SGD require compared to Adam?
    \item What is your personal workflow for optimizer selection?
\end{itemize}

\textit{Red flags to avoid:}
\begin{itemize}
    \item Having no opinion or saying ``always use Adam''
    \item Not mentioning the tuning requirement difference
    \item Forgetting that SGD needs momentum in practice
\end{itemize}

\textit{Strong answer key points:}
\begin{itemize}
    \item \textbf{Choose SGD when}: (1) Training CNNs for image classification where final accuracy matters, (2) You have compute budget for hyperparameter search, (3) Generalization is more important than training speed
    \item \textbf{Choose Adam when}: (1) Training Transformers/LLMs (Adam is standard), (2) Rapid prototyping, (3) RNNs or problems with sparse gradients, (4) Limited tuning budget
    \item \textbf{Practical workflow}: Start with AdamW (default $\eta=0.001$). If training works but generalization is poor, try SGD+momentum ($\eta=0.1$, $\gamma=0.9$) with learning rate search
    \item \textbf{Domain-specific}: Vision: SGD often wins. NLP: Adam/AdamW. RL: Adam. GANs: Adam with $\beta_1=0.5$
\end{itemize}

\vspace{0.5em}
\textbf{Question 5: What is learning rate warmup and why does it help?}

\textit{Expected follow-ups:}
\begin{itemize}
    \item How long should warmup be?
    \item Why is warmup especially important for Transformers?
    \item What happens if you skip warmup with a large learning rate?
\end{itemize}

\textit{Red flags to avoid:}
\begin{itemize}
    \item Not knowing that warmup is critical for large models and batch sizes
    \item Confusing warmup with learning rate decay
    \item Thinking warmup is only about numerical stability
\end{itemize}

\textit{Strong answer key points:}
\begin{itemize}
    \item \textbf{What it is}: Gradually increase learning rate from near-zero to target value over first $W$ steps
    \item \textbf{Why it helps}: (1) Adam's moment estimates are biased toward zero initially, need time to stabilize; (2) Large early updates can permanently damage representations; (3) Initial gradients may not represent steady-state statistics
    \item \textbf{For Transformers}: Attention layers are particularly sensitive; without warmup, attention patterns can lock into degenerate configurations early
    \item \textbf{Typical duration}: 1-10\% of total training; for LLMs often 1000-4000 steps
    \item \textbf{For large batch training}: Linear scaling rule requires warmup---jumping to high LR without warmup causes instability
    \item \textbf{Common schedule}: Linear warmup followed by cosine decay to near-zero
\end{itemize}

\end{interviewquestions}

%------------------------------------------------------------------------------
%  CHAPTER SUMMARY
%------------------------------------------------------------------------------
\begin{chaptersummary}
\textbf{What We Covered:}
\begin{itemize}
    \item Adaptive methods (AdaGrad, RMSprop) automatically scale per-parameter learning rates based on gradient history
    \item Adam combines momentum (Chapter~\ref{ch:gradient-descent}) with adaptive learning rates; bias correction is crucial for stable early training
    \item AdamW decouples weight decay from gradient updates for better generalization
    \item Learning rate schedules (warmup, cosine annealing, warm restarts) adapt optimization over training
\end{itemize}

\textbf{Key Formulas:}
\begin{align*}
    &\text{Adam: } \vect{m}_t = \beta_1 \vect{m}_{t-1} + (1-\beta_1)\vect{g}_t, \quad \vect{v}_t = \beta_2 \vect{v}_{t-1} + (1-\beta_2)\vect{g}_t^2 \\
    &\text{Bias Correction: } \hat{\vect{m}}_t = \frac{\vect{m}_t}{1-\beta_1^t}, \quad \hat{\vect{v}}_t = \frac{\vect{v}_t}{1-\beta_2^t} \\
    &\text{Update: } \vect{\theta}_{t+1} = \vect{\theta}_t - \eta \frac{\hat{\vect{m}}_t}{\sqrt{\hat{\vect{v}}_t} + \epsilon}
\end{align*}

\textbf{When to Use This:}
\begin{itemize}
    \item AdamW with warmup for Transformer/LLM training (essential combination)
    \item SGD+Momentum for CNNs when maximizing final accuracy matters
    \item Adam for quick prototyping---robust defaults, minimal tuning required
    \item RMSprop for RNNs or when gradient magnitudes vary significantly across time
\end{itemize}

\textbf{Next Up:} Regularization techniques that prevent overfitting, from explicit penalties (L1/L2) to implicit methods (dropout, early stopping).
\end{chaptersummary}

\end{document}
